\newpage
\section{\textbf{KẾT LUẬN}}

Thông qua quá trình đọc mã nguồn, đối chiếu tài liệu và tổ chức lại các artifact kiểm thử của hệ thống Baby Bliss, báo cáo đã xác định được một phạm vi nghiên cứu rõ ràng cho môn Kiểm thử phần mềm. Trong phạm vi đó, nhóm lựa chọn 5 chức năng tiêu biểu gồm quản lý mã khuyến mãi, thanh toán -- giỏ hàng, yêu thích, profile và đăng nhập -- xác thực để triển khai các nội dung phân tích. Tuy nhiên, thay vì trải đều cho tất cả chức năng, báo cáo tập trung sâu hơn vào những phần có nhiều bằng chứng hơn trong mã nguồn và trong bộ test, nhờ đó giữ được tính thực tế và sự nhất quán giữa các chương.

Kết quả quan trọng nhất của báo cáo là đã tạo được một mạch liên kết tương đối rõ từ xác định phạm vi chức năng $\rightarrow$ thiết kế test case black box $\rightarrow$ phân tích white box $\rightarrow$ đối chiếu với unit test $\rightarrow$ liên hệ với phần hiện thực và demo kiểm thử tự động. Mạch này thể hiện rõ nhất ở quản lý mã khuyến mãi và đăng nhập. Trong đó, quản lý mã khuyến mãi là chức năng hoàn thiện nhất vì được theo dõi xuyên suốt từ bảng quyết định, phân tích hàm \code{validateCoupon} đến kết quả unit test; đăng nhập là chức năng có mức độ hoàn thiện tiếp theo với sự nối tiếp khá rõ từ phân vùng tương đương, phủ nhánh cho hàm \code{login} đến phần thực thi kiểm thử. Đối với profile, báo cáo đã thể hiện được các ràng buộc dữ liệu và có liên hệ với demo, nhưng mức độ hoàn thiện vẫn thấp hơn vì chưa có một mục unit test riêng đủ mạnh như hai chức năng trên.

Một điểm đáng ghi nhận khác là báo cáo đã đưa được các bằng chứng thực thi vào nội dung thay vì chỉ dừng ở mức đề xuất test case. Các kết quả chạy Jest cho quản lý mã khuyến mãi, đăng nhập và inventory cho thấy một phần quan trọng của phân tích đã được chuyển thành test có thể chạy được, trong khi chương hiện thực và demo tiếp tục làm rõ mối liên hệ giữa backend, frontend, unit test và Playwright. Dù vậy, độ sâu kiểm thử giữa các chức năng vẫn chưa đồng đều: yêu thích hiện mạnh hơn ở góc độ black box và mô hình trạng thái, thanh toán -- giỏ hàng chủ yếu mới dừng ở mức xác định phạm vi, còn inventory chủ yếu đóng vai trò minh chứng mở rộng về white box và thực thi chứ chưa phải là một trục phân tích xuyên suốt từ đầu.

Từ các kết quả trên, hướng hoàn thiện tiếp theo của báo cáo là tiếp tục cân bằng độ sâu giữa các chương và giữ chặt mối liên hệ giữa phần phân tích với phần thực thi. Cụ thể, nhóm có thể mở rộng thêm chiều sâu cho chức năng yêu thích, bổ sung bằng chứng unit test cho profile, làm dày hơn phần thanh toán -- giỏ hàng nếu tiếp tục giữ chức năng này trong phạm vi báo cáo, đồng thời hoàn thiện thêm các artifact kiểm thử để việc đối chiếu khi bảo vệ được rõ ràng hơn. Nếu tiếp tục phát triển theo hướng đó, báo cáo sẽ không chỉ đáp ứng yêu cầu môn học mà còn thể hiện rõ hơn cách lựa chọn kỹ thuật kiểm thử phù hợp và cách biến kết quả phân tích thành các bằng chứng thực thi trên hệ thống thực.
